\section{Getting Started}
\label{sec:GettingStarted}

This section walks through bringing up a three-node Raft cluster and issuing
the first commands. It is intended as an orientation for both application
developers evaluating the library and machinery developers who need a working
cluster to test against.

All three implementations (Java, C++, Rust) are covered. Pick the one that
matches your build toolchain, or mix them using the Docker Compose files
described in \S\ref{sec:GettingStartedDocker}.

\subsection{Building from source}

\subsubsection{Java}

\begin{lstlisting}[language=bash,caption={Java build}]
# Build all Java modules, run unit tests, produce raft-dist/target/raft-*.jar
mvn -q package
\end{lstlisting}

The distributable JAR lives at \filepath{raft-dist/target/raft-1.0-SNAPSHOT.jar}.
It bundles all runtime dependencies and can be invoked directly:

\begin{lstlisting}[language=bash,caption={Java CLI}]
java -jar raft-dist/target/raft-1.0-SNAPSHOT.jar n1@127.0.0.1:11081
\end{lstlisting}

\subsubsection{C++}

\begin{lstlisting}[language=bash,caption={C++ build}]
cmake -S graft-cpp -B graft-cpp/build
cmake --build graft-cpp/build
\end{lstlisting}

Produces \filepath{graft-cpp/build/graft\_smoke}, a single binary covering
all server, client and probe commands.

\subsubsection{Rust}

\begin{lstlisting}[language=bash,caption={Rust build}]
cd graft-rust
cargo build -p graft-app-kv
\end{lstlisting}

Produces \filepath{graft-rust/target/debug/graft-kv} (or
\filepath{target/release/graft-kv} with \texttt{--release}). The binary
provides 26 subcommands covering server, client, operational and smoke modes.

\subsection{Starting a local cluster}

The simplest way to start a cluster is with the active persistent server mode.
Each node needs a unique port and a peer list containing all other nodes.

\subsubsection{Java cluster}

Start three terminals (or background processes):

\begin{lstlisting}[language=bash,caption={Java 3-node cluster}]
java -jar raft-dist/target/raft-1.0-SNAPSHOT.jar \
  n1@127.0.0.1:11081 n2@127.0.0.1:11082 n3@127.0.0.1:11083

java -jar raft-dist/target/raft-1.0-SNAPSHOT.jar \
  n2@127.0.0.1:11082 n1@127.0.0.1:11081 n3@127.0.0.1:11083

java -jar raft-dist/target/raft-1.0-SNAPSHOT.jar \
  n3@127.0.0.1:11083 n1@127.0.0.1:11081 n2@127.0.0.1:11082
\end{lstlisting}

\subsubsection{Rust cluster}

The Rust CLI uses explicit \texttt{serve-active-persistent} with positional
arguments instead of a combined peer-spec/as-argument convention:

\begin{lstlisting}[language=bash,caption={Rust 3-node cluster}]
graft-rust/target/debug/graft-kv serve-active-persistent \
  127.0.0.1 11081 n1 /tmp/n1.state 0 0 0 \
  n2@127.0.0.1:11082 n3@127.0.0.1:11083

graft-rust/target/debug/graft-kv serve-active-persistent \
  127.0.0.1 11082 n2 /tmp/n2.state 0 0 0 \
  n1@127.0.0.1:11081 n3@127.0.0.1:11083

graft-rust/target/debug/graft-kv serve-active-persistent \
  127.0.0.1 11083 n3 /tmp/n3.state 0 0 0 \
  n1@127.0.0.1:11081 n2@127.0.0.1:11082
\end{lstlisting}

\subsubsection{Verifying cluster health}

Once all nodes are started, one will be elected leader within a few seconds.
Verify with the telemetry or cluster-summary commands (both work against any
implementation since they share the same wire protocol):

\begin{lstlisting}[language=bash,caption={Cluster health check (C++ probe)}]
graft-cpp/build/graft_smoke cluster-summary 127.0.0.1 11081 cpp-cli
graft-cpp/build/graft_smoke telemetry 127.0.0.1 11082 true
\end{lstlisting}

Or using the Rust binary:

\begin{lstlisting}[language=bash,caption={Cluster health check (Rust probe)}]
graft-rust/target/debug/graft-kv cluster-summary 127.0.0.1 11081
graft-rust/target/debug/graft-kv telemetry 127.0.0.1 11082 \
  --require-leader-summary
\end{lstlisting}

Look for \texttt{state: LEADER} on one node and \texttt{state: FOLLOWER} on
the others.

\subsection{First commands}

Put a value and read it back. The client must target the leader; followers
redirect.

\begin{lstlisting}[language=bash,caption={Put and get (C++ probe)}]
graft-cpp/build/graft_smoke client-put 127.0.0.1 11081 greeting hello my-client
graft-cpp/build/graft_smoke client-get 127.0.0.1 11081 greeting my-client
\end{lstlisting}

Or with Rust:

\begin{lstlisting}[language=bash,caption={Put and get (Rust probe)}]
graft-rust/target/debug/graft-kv client-put 127.0.0.1 11081 \
  greeting hello my-client
graft-rust/target/debug/graft-kv client-get 127.0.0.1 11081 \
  greeting my-client
\end{lstlisting}

If you send a client-get to a follower, the response will include
\texttt{status: REDIRECT} with the leader's address.

\subsection{Docker deployment}
\label{sec:GettingStartedDocker}

For multi-language clusters or reproducible testing, use Docker Compose.
Four Compose files are provided:

\begin{itemize}
\item \filepath{docker/compose-srv-java.yml} --- 3 Java nodes and 1 CoreDNS
\item \filepath{docker/compose-srv-cpp.yml} --- 3 C++ nodes and 1 CoreDNS
\item \filepath{docker/compose-srv-rust.yml} --- 3 Rust nodes and 1 CoreDNS
\item \filepath{docker/compose-srv-mixed.yml} --- 1 Java, 1 C++, 1 Rust and 1 CoreDNS
\end{itemize}

Start a cluster with:

\begin{lstlisting}[language=bash,caption={Docker smoke test}]
docker/run-srv-smoke.sh java          # or cpp, rust, mixed
docker/run-srv-suite.sh               # all four in sequence
\end{lstlisting}

The smoke script builds images, starts containers, waits for leader election,
probes all three nodes, and tears down on exit. All nodes share a CoreDNS
container resolving the \texttt{raft.local} zone. Host ports 17001--17003
map to container port 7000.

\subsection{DNS SRV bootstrapping}
\label{sec:DnsSrvBootstrapping}

The local cluster examples above use explicit peer lists
(\texttt{n2@127.0.0.1:11082,n3@...}) so each node knows its peers at startup.
In containerized deployments---Docker Compose, Kubernetes, OpenShift---this
becomes unwieldy. DNS SRV records let nodes discover peers from a single
service name instead of a static list.

\subsubsection{Why DNS SRV?}

A DNS SRV record maps a service name to multiple \emph{target} entries, each
specifying a hostname and port. This matches the problem of cluster
bootstrapping: a new node needs to contact \emph{any} existing member, and
the DNS zone is the natural place to publish the current member set.

SRV records are used natively in Kubernetes and OpenShift via headless
Services. A StatefulSet-backed Raft cluster can publish SRV records
automatically, giving each node a stable identity \\
(\texttt{raft-0.raft.\textless ns\textgreater .svc.cluster.local}) and a
discoverable service endpoint \\ 
(\texttt{\_raft.\_tcp.raft.\textless ns\textgreater .svc.cluster.local}). The same mechanism works in Docker
Compose with a CoreDNS sidecar.

The startup contract is intentionally minimal. A node needs only three
pieces of information:

\begin{itemize}
\item \texttt{RAFT\_NODE\_ID} --- its logical Raft identity,
\item \texttt{RAFT\_BIND\_HOST}/\texttt{RAFT\_BIND\_PORT} --- its listen address,
\item \texttt{RAFT\_CLUSTER\_SRV} --- the SRV service name to resolve.
\end{itemize}

Everything else---current membership, leader identity, cluster metadata---is
learned from the first contacted peer via the Raft protocol itself.

\subsubsection{SRV record format}

An SRV record has the form:

\begin{lstlisting}[caption={SRV record format}]
_service._proto.name. TTL IN SRV priority weight port target.
\end{lstlisting}

For a three-node Raft cluster in the \texttt{raft.local} zone:

\begin{lstlisting}[caption={Example SRV zone entries}]
_raft._tcp.raft.local. IN SRV 10 10 7000 raft-1.raft.local.
_raft._tcp.raft.local. IN SRV 10 10 7000 raft-2.raft.local.
_raft._tcp.raft.local. IN SRV 10 10 7000 raft-3.raft.local.
\end{lstlisting}

All three entries have equal priority (10) and weight (10) since all Raft
voters are equally valid bootstrap contacts. The port (7000) is the Raft
TCP port each node listens on.

\subsubsection{Startup sequence}

When a node starts with \texttt{RAFT\_CLUSTER\_SRV} set, it executes the
following sequence:

\begin{enumerate}
\item Bind the local Raft port.
\item Resolve \texttt{RAFT\_CLUSTER\_SRV} via DNS (using libresolv in C++,
JNDI in Java, system \texttt{dig} in Rust).
\item Shuffle the returned SRV targets to avoid all nodes contacting the
same seed.
\item Connect to each target in turn until one responds with a valid Raft
handshake.
\item From that contacted peer, learn: the current term, the current leader,
and the committed cluster membership.
\item Stop treating DNS as authoritative---subsequent peer discovery is
driven by Raft membership log entries, not repeated SRV lookups.
\end{enumerate}

If no seed answers and \texttt{RAFT\_BOOTSTRAP\_NEW\_CLUSTER} is not set,
the node fails fast. This prevents a transient DNS failure from
accidentally forming a split-brain cluster. Bootstrapping a new cluster
must be explicit.

\subsubsection{Implementation across languages}

All three implementations support DNS SRV bootstrapping, though the
resolution mechanism differs:

\begin{description}
\item[Java] \texttt{DnsSrvSeedProvider} resolves via JNDI DNS lookups
(\texttt{javax.naming.directory}). Set \texttt{RAFT\_CLUSTER\_SRV} as a
system property or environment variable. The Java adapter reads it from
\texttt{RuntimeConfiguration}.
\item[C++] \texttt{DnsSrvSeedProvider} uses the system resolver
(\texttt{res\_query} from libresolv). The C++ CLI reads
\texttt{RAFT\_CLUSTER\_SRV} from the environment. The CMake build links
\texttt{resolv} on Unix.
\item[Rust] \texttt{DnsSrvSeedProvider} shells out to \texttt{dig +short SRV}
with a fallback to \texttt{nslookup -type=SRV}. The Docker entrypoint
(\filepath{entrypoint-rust.sh}) reads \texttt{RAFT\_CLUSTER\_SRV} and calls
\texttt{dig} to build the peer list before launching \texttt{graft-kv}.
\end{description}

In Docker Compose, the Rust entrypoint is used with explicit
\texttt{RAFT\_PEERS} by default, which avoids the \texttt{dig} dependency
in the container. Adding \texttt{dnsutils} to the Rust Dockerfile enables
the SRV path. The Java and C++ images support SRV natively without extra
packages.

\subsubsection{Kubernetes / OpenShift alignment}

The DNS SRV model maps directly to Kubernetes headless Services:

\begin{lstlisting}[style=yaml, caption={Kubernetes headless Service}]
apiVersion: v1
kind: Service
metadata:
  name: raft
spec:
  clusterIP: None        # headless
  selector:
    app: raft
  ports:
    - name: raft
      protocol: TCP
      port: 7000
      targetPort: 7000
\end{lstlisting}

With a StatefulSet backing this Service, Kubernetes creates SRV records
automatically. Each node sets
\texttt{RAFT\_CLUSTER\_SRV=\_raft.\_tcp.raft.\textless ns\textgreater .svc.cluster.local}
and the platform handles DNS propagation. This gives a clean portability
boundary: locally you own the CoreDNS zone file; in OpenShift, the
platform owns the DNS records.

\subsection{Where to go next}

\begin{itemize}
\item Section~\ref{sec:ApplicationDeveloper} for the domain application boundary.
\item Section~\ref{sec:ConsistencyBoundaries} for what Raft guarantees---and does
not guarantee---under partitions and failovers.
\item Section~\ref{sec:RaftAndXatmi} for a comparison with XATMI/XA transaction
models, useful if you are coming from a traditional middleware background.
\item Section~\ref{sec:JavaLibrary}, \ref{sec:CppLibrary}, or \ref{sec:RustLibrary}
for language-specific integration details.
\item Section~\ref{sec:TestingAndOperations} for test workflows and mixed-smoke checks.
\item \filepath{raft-app-kv/}, \filepath{graft-cpp/src/app/}, or
\filepath{graft-rust/graft-app-kv/} for reference implementations.
\end{itemize}
